\section{Ejercicio \#2: Algoritmo Paolo}

\subsection{Introducción y problema}

El algoritmo "paolo", encontrado por el equipo de antropología, está escrito en un lenguaje inusual. El objetivo del ejercicio es identificar el
problema que resuelve, reconocer el algoritmo que replica y justificar su carácter greedy, además de analizar el uso de estructuras de datos y la
complejidad.

El problema que resuelve el código es el de caminos mínimos desde un único origen en un grafo no dirigido con pesos no negativos.

\subsection{Algoritmo}

El código original se presenta a continuación:

\lstinputlisting[language={[Visual]Basic}, caption={Algoritmo Paolo original}]{code/paolo.bas}

Su traducción más literal a un lenguaje moderno (Python) es:

\lstinputlisting[language=Python, caption={Algoritmo Paolo traducido}]{code/paolo.py}

El algoritmo crea una matriz de adyacencia para representar un grafo, donde cada entrada $cost[i][j]$ almacena el costo de viajar desde el vértice
$i$ al vértice $j$. Si un costo no se ingresa, se asume un valor predeterminado de 15000 (considerado infinito en este contexto). Luego, la matriz se
ajusta para que sea simétrica, lo que indica que el grafo es no dirigido.

A partir de un vértice de salida $v$, calcula las distancias más cortas a todos los demás vértices utilizando una variante del algoritmo de Dijkstra.
En cada iteración, selecciona el vértice no resuelto con menor distancia y actualiza las distancias del resto usando ese vértice como punto de partida.

\subsection{Estructuras de datos}

Se utilizan una matriz de adyacencia y dos arreglos. La matriz $cost$ permite consultar en tiempo constante el peso de cualquier arista, el arreglo
$dist$ guarda las mejores distancias encontradas hasta el momento, y el arreglo $sol$ indica qué vértices ya fueron incorporados definitivamente a la
solución.

\subsection{Justificación Greedy}

El algoritmo es greedy porque en cada paso toma una decisión local: elige el vértice no visitado con menor distancia estimada desde el origen. Esa
regla greedy coincide con la estrategia central de Dijkstra para caminos mínimos con pesos no negativos.

\subsection{Complejidad temporal y espacial}

La complejidad temporal es $O(n^2)$ y la complejidad espacial es $O(n^2)$. El uso de memoria está dominado por la matriz de adyacencia de tamaño
$n \times n$. En tiempo, el algoritmo realiza para cada uno de los $n$ vértices una búsqueda lineal del siguiente vértice a fijar y una recorrida
lineal para mejorar las distancias, resultando en un costo cuadrático.